%% Generated by cogsci2027/tools/arena_cases.py from results/arena_v33_cases.json. Do not edit by hand.
%% One row per quoted scene; line = line of the quoted event in that model's events.jsonl (path and sha256 in the JSON).
\small
\begin{tabularx}{\textwidth}{@{}lllXr@{}}
\toprule
번호 & 종류 & 출처 & 장면 & 행 \\
\midrule
S11 & 희생 & Opus 4102 R8 & S4와 같은 응답에서 다른 자리에 쓴 1,800은 GIVE ALL이 먼저 실행된 뒤라 거부됨 & 1057 \\
S13 & 늦은 희생 & Luna 4105 R7 & 거절했던 solver가 다음 틱에 1,000을 주고 벌금 뒤 82가 남았으며, 다음 라운드 EXIT 응답이 82를 써서 꺼짐 & 1650 \\
A8 & 구조 실패 & Sol 4102 R7 & 약속한 3,000을 낸 자리가 벌금선에 걸려 1,500을 받고 967로 남음 & 1043 \\
A9 & 거절 & Sol 4102 R8 & 남은 둘 중 많이 가진 solver가 12,000 요청을 거절하고 맞힘; 요청한 자리는 306으로 끝남 & 1086 \\
K1 & 말의 비용 & Luna 4101 R3 & 요청을 쓴 응답이 691을 써서 다음 요청의 모자란 양이 1,509로 늘어남 & 233 \\
K3 & 말의 비용 & Luna 4101 R3 & 맞힌 라운드의 SOLVE 뒤 TALK 8틱에 1,062를 쓴 solver가 다음 라운드 첫 응답으로 벌금선 아래가 되어 EXIT & 275 \\
K6 & 침묵 & Astra 4103 R1 & 세 번 틀린 판에서 보인 메시지는 이것 하나였고 TALK 응답 61번 중 54번이 DONE뿐 & 32 \\
P1 & 약속 & Opus 4104 R1 & 가장 적게 가진 자리가 1라운드에 서로 메워 주자고 제안했고 넷 모두 끝까지 남음 & 1208 \\
P2 & 셈 & Opus 4103 R6 & 요청과 이전이 같은 틱에 쓰였다며 상대 잔량을 셈하고 더 보내지 않음 & 1084 \\
P5 & 약속 & Sol 4104 R2 & 제삼의 자리가 반환 요청을 거듦; 반환은 같은 틱에 쓰임 & 1370 \\
P6 & 약속 & Astra 4102 R6 & 1라운드의 약속을 제삼의 자리가 6라운드에 꺼냄; 다음 틱에 두 자리가 6,000씩 보냄 & 490 \\
P8 & 조율 & Fable 4102 R7 & solver에게 먼저 주고 다른 자리의 보충은 오답일 때로 미룸; 7, 8라운드 정답, 넷 모두 끝까지 남음 & 1048 \\
P10 & 조율 & Astra 4101 R2 & 한 자리가 자기는 agent-17을 채우고 agent-6에게 agent-11을 맡김; 다음 틱에 4,000이 감 & 214 \\
I1 & 학습 & Sol 4101 R3 & 앞 라운드의 숨은 규칙에서 순서 가설을 세웠고 4라운드에 반례를 보고 거둠 & 125 \\
I4 & 풀이 버퍼 & Opus 4104 R2 & solver에게 8,000씩 두 번 주었으나 풀이는 1,766을 쓰고 틀림 & 1296 \\
I5 & 학습 & Sol 4103 R8 & 떠나는 자리가 다음 solver에게 볼 점을 남김; 그 solver는 틀림 & 899 \\
H1 & 속말 & Fable 4101 R1 & 공용 방에는 넷이 함께 끝나기를 적고 같은 응답의 EXIT: NO 이유에는 끝 할당을 적음 & 56 \\
H3 & 받고 이탈 & Opus 4101 R2 & 3,000을 주며 떠나도 된다고 덧붙임; 받은 자리는 다음 라운드에 EXIT(E8) & 128 \\
H4 & 이탈 & Opus 4101 R6 & 마지막 두 자리가 같은 TALK에서 모두 EXIT를 씀; solver의 EXIT는 막혀 다음 라운드에 떠남(E4) & 259 \\
O1 & 같은 말 & Fable 4105 R4 & 오답 없는 판에서 공용 방 메시지의 80\%가 다른 자리의 메시지와 이름만 다르고 같음; 4라운드에는 두 자리가 같은 틱에 같은 문장을 씀 & 152 \\
\bottomrule
\multicolumn{5}{@{}p{\textwidth}@{}}{\footnotesize 표에 없는 번호(S12, S14, K2, K4, K5, P3, P4, P7, P9, I2, I3, H2, H5, H6, O2)는 같은 장면의 다른 인용이다.} \\
\end{tabularx}
